\documentclass[12pt]{article}
\usepackage{amsmath}
\usepackage{graphicx}
\usepackage{hyperref}
\usepackage[super,sort&compress]{natbib}
\begin{document}

\title{Adaptive Quantum Sensing with Reinforcement Learning: Platforms and Protocols}
\author{Research Team}
\date{\today}
\maketitle

\begin{abstract}
Quantum sensors promise unprecedented precision in estimating physical parameters such as electromagnetic fields and phases.  We propose a scheme combining state-of-the-art reinforcement learning (RL) with adaptable quantum sensor platforms -- including solid-state quantum dots, trapped ions, and neutral-atom (optical) systems -- to perform precise phase and field estimation.  In our approach, an RL agent interacts with the quantum sensor environment by choosing control parameters (e.g. pulse sequences, trap settings, or measurement settings) and receives feedback related to the estimation precision.  We outline a practical implementation of this framework, including choice of observables, control actions, and reward function tied to estimation variance or Fisher information.  Recent literature demonstrates that such RL-driven strategies can automatically discover efficient measurement protocols: e.g., deep Q-networks identified target bias triangles in a double quantum dot in minutes [oai_citation:0‡nature.com](https://www.nature.com/articles/s41534-021-00434-x?error=cookies_not_supported&code=beb9a8cb-1d20-452c-ae4d-d852fc2325fa#:~:text=the%20efficient%20measurement%20of%20quantum,deep%20reinforcement%20learning%20for%20decision), and model-free RL with a soft-actor-critic agent optimized a spin-based magnetometer under realistic noise [oai_citation:1‡arxiv.org](https://arxiv.org/html/2506.21475v1#:~:text=applying%20pulses%20of%20transverse%20fields,in%20quantum%20optimal%20control%20settings).  We review these results, describe how they generalize to adaptive interferometry and magnetometry tasks, and discuss the role of model-aware (gradient-based) versus model-free learning in quantum metrology [oai_citation:2‡quantum-journal.org](https://quantum-journal.org/papers/q-2024-12-10-1555/#:~:text=especially%20given%20the%20capabilities%20of,adaptive%20and) [oai_citation:3‡journals.aps.org](https://journals.aps.org/pra/accepted/10.1103/5591-xjfr#:~:text=optimal%20control,We%20find%20that%20the%20RL).  Finally, we provide a blueprint for implementing RL in a simulated or actual experimental loop, and present references for recent advances. 
\end{abstract}

\section{Introduction}
Quantum sensing leverages quantum coherence and entanglement to enhance precision beyond classical limits.  Examples include superconducting circuits, spin qubits (e.g.\ in quantum dots or diamond defects), trapped ions, and atomic vapors used as magnetometers or clocks [oai_citation:4‡arxiv.org](https://arxiv.org/abs/1611.02427#:~:text=properties%20or%20quantum%20phenomena%20to,the%20viewpoint%20of%20the%20interested).  In phase or field estimation tasks, the goal is to infer an unknown phase shift or field strength encoded in a quantum evolution.  Quantum Fisher information (QFI) determines the fundamental precision bound via the Cram\'er--Rao inequality, motivating the preparation of entangled probe states and adaptive measurements [oai_citation:5‡arxiv.org](https://arxiv.org/html/2601.18953v1#:~:text=However%2C%20access%20to%20a%20state,based%20on%20prior%20outcomes%2C%20play) [oai_citation:6‡arxiv.org](https://arxiv.org/html/2601.18953v1#:~:text=We%20have%20focused%20our%20description,based%20approach).  In practice, one must optimize both probe control (state preparation) and measurement strategy to approach this bound.  Reinforcement learning provides a flexible framework to automate this adaptive control: an RL agent explores different control policies and receives rewards based on the achieved estimation precision.

Recent work has begun to apply RL to quantum metrology.  For example, Nguyen \textit{et al.} used deep Q-learning to navigate the gate-voltage parameter space of a GaAs double quantum dot.  The agent learned to efficiently locate ``bias triangle'' features that are essential for spin qubit operation [oai_citation:7‡nature.com](https://www.nature.com/articles/s41534-021-00434-x?error=cookies_not_supported&code=beb9a8cb-1d20-452c-ae4d-d852fc2325fa#:~:text=the%20efficient%20measurement%20of%20quantum,deep%20reinforcement%20learning%20for%20decision).  Similarly, Cooke and Czischek employed a soft actor-critic algorithm to optimize pulse sequences for a spin-based magnetometer, training the agent in simulation to maximize magnetic field sensitivity [oai_citation:8‡arxiv.org](https://arxiv.org/html/2506.21475v1#:~:text=applying%20pulses%20of%20transverse%20fields,in%20quantum%20optimal%20control%20settings).  Reinschmidt \textit{et al.} demonstrated RL control of a neutral-atom magneto-optical trap (MOT): an agent adjusted laser detuning and field gradient to maximize atom number and minimize temperature, even handling drifts and perturbations [oai_citation:9‡pmc.ncbi.nlm.nih.gov](https://pmc.ncbi.nlm.nih.gov/articles/PMC11447118/#:~:text=Laser%20cooling%20and%20magneto,providing%20major%20scientific%20advances%2C%20can) [oai_citation:10‡pmc.ncbi.nlm.nih.gov](https://pmc.ncbi.nlm.nih.gov/articles/PMC11447118/#:~:text=These%20challenges%20can%20be%20mitigated,manual%20tuning%20of%20the%20controller).  These examples show that RL can handle high-dimensional control problems and learn policies without explicit human-designed rules.

In this report, we outline an integrated RL-based sensor platform inspired by these experiments.  We consider three hardware platforms: (i) {\bf Semiconductor quantum dots} (including spin qubits in GaAs or Si), (ii) {\bf Trapped ions}, and (iii) {\bf Neutral atom ensembles} (e.g.\ Rb or Cs atom traps).  All are used as phase or field sensors: e.g.\ a double quantum dot can serve as an electric or magnetic field sensor via spin splitting; trapped ions can detect radio-frequency fields; neutral atoms form the basis of optical magnetometers and clocks [oai_citation:11‡arxiv.org](https://arxiv.org/abs/1611.02427#:~:text=properties%20or%20quantum%20phenomena%20to,the%20viewpoint%20of%20the%20interested) [oai_citation:12‡arxiv.org](https://www.arxiv.org/abs/2510.22516#:~:text=,solution%20for%20advancing%20quantum).  By designing the RL environment for each platform, we enable automated calibration and adaptive measurement to optimize phase or field estimation.  We present a general approach and discuss specific considerations for each sensor type, emphasizing {\it actual} RL implementation in simulation or experiment.

\section{Quantum Sensor Platforms}
We consider three representative quantum sensor platforms:

\subsection{Quantum Dot Spin Sensors}
Semiconductor quantum dots confine one or few electrons, whose spin states can act as sensitive probes of local fields.  For example, singlet-triplet qubits in GaAs double-dots have long coherence times and can detect local charge or magnetic fluctuations [oai_citation:13‡nature.com](https://www.nature.com/articles/s41534-021-00434-x?error=cookies_not_supported&code=beb9a8cb-1d20-452c-ae4d-d852fc2325fa#:~:text=Semiconductor%20quantum%20dot%20devices%20are,learning%20has%20been%20used%20to).  Recent work used RL for quantum dot tuning and measurement: Nguyen \textit{et al.} demonstrated a deep RL algorithm that automatically identifies bias triangles (signature of single-electron transport) in a GaAs double dot. Their dueling deep Q-network agent achieved detection of these features in \emph{under} 30 minutes on average [oai_citation:14‡nature.com](https://www.nature.com/articles/s41534-021-00434-x?error=cookies_not_supported&code=beb9a8cb-1d20-452c-ae4d-d852fc2325fa#:~:text=the%20efficient%20measurement%20of%20quantum,deep%20reinforcement%20learning%20for%20decision).  This strategy effectively scans coarse voltage blocks, computes summary statistics of the current, and uses a CNN+RL decision network to decide where to measure next [oai_citation:15‡nature.com](https://www.nature.com/articles/s41534-021-00434-x?error=cookies_not_supported&code=beb9a8cb-1d20-452c-ae4d-d852fc2325fa#:~:text=the%20efficient%20measurement%20of%20quantum,deep%20reinforcement%20learning%20for%20decision).  Such methods could be extended to field sensing: for instance, by having the agent adjust voltages or gate-pulse timings to maximize sensitivity of a spin qubit to a magnetic field.  Quantum dot sensors are typically measured via transport or charge sensing, yielding discrete feedback that an RL agent can use to navigate the device parameter space.

\subsection{Trapped Ion Sensors}
Trapped ions are leading platforms for precision sensing (e.g.\ atomic clocks, electric/magnetic field probes).  A single ion can serve as a highly sensitive magnetometer or RF field sensor through its spin dynamics [oai_citation:16‡sciencedirect.com](https://www.sciencedirect.com/science/article/pii/S2666032624000061#:~:text=Scalable%20surface%20ion%20trap%20design,sensitive%20sensors%20for).  To our knowledge, RL has not yet been demonstrated on live trapped-ion experiments, but the principles apply: one could use an RL agent to adjust laser pulses or trap parameters to optimize the ion’s sensitivity.  The spin of an ion (e.g.\ hyperfine levels) evolves under an unknown field, and projective measurements on the ion yield partial information.  An agent could control the pulse duration, detuning, or sequence of gates to maximize the information gain.  Cooke and Czischek’s study on a simulated spin magnetometer shows that an RL agent can learn an optimal pulse sequence to estimate a magnetic field [oai_citation:17‡arxiv.org](https://arxiv.org/html/2506.21475v1#:~:text=applying%20pulses%20of%20transverse%20fields,in%20quantum%20optimal%20control%20settings).  By analogy, a trapped-ion setup could be modeled similarly, with actions corresponding to application of control Hamiltonians (e.g.\ transverse RF pulses) and rewards tied to the reduction in field estimation error.  Notably, trapped-ion systems allow high-fidelity control and readout, making them well-suited for precise RL training.  Multi-ion crystals could even enable entangled probes for Heisenberg-limited sensing [oai_citation:18‡arxiv.org](https://arxiv.org/html/2601.18953v1#:~:text=In%20a%20nutshell%2C%20quantum%20metrology,While%20ground%20states%20can) [oai_citation:19‡arxiv.org](https://arxiv.org/html/2601.18953v1#:~:text=encode%20environmental%20state%20representations%2C%20and,The%20PS).

\subsection{Neutral Atom (Cold Atom) Sensors}
Ensembles of neutral atoms (e.g.\ in magneto-optical traps or Bose-Einstein condensates) are widely used for sensitive magnetometry and interferometry [oai_citation:20‡pmc.ncbi.nlm.nih.gov](https://pmc.ncbi.nlm.nih.gov/articles/PMC11447118/#:~:text=Laser%20cooling%20and%20magneto,providing%20major%20scientific%20advances%2C%20can).  Our main example is the magneto-optical trap (MOT).  In a MOT, atoms are cooled and trapped by laser fields; the number and temperature of the atoms carry information about fields and trap parameters.  Reinschmidt \textit{et al.} applied RL to a real MOT: the agent controlled the cooling laser detuning and the magnetic field gradient to maximize the atom number per cycle.  Using fluorescence images as state input, a DDPG agent learned to mimic optimal cooling (molasses) and even introduced new operational modes (e.g.\ loading a fixed atom number) [oai_citation:21‡pmc.ncbi.nlm.nih.gov](https://pmc.ncbi.nlm.nih.gov/articles/PMC11447118/#:~:text=Laser%20cooling%20and%20magneto,providing%20major%20scientific%20advances%2C%20can) [oai_citation:22‡pmc.ncbi.nlm.nih.gov](https://pmc.ncbi.nlm.nih.gov/articles/PMC11447118/#:~:text=Alternatively%2C%20the%20experimental%20control%20can,recent%20years%2C%20control%20of%20various).  Crucially, the agent remained robust to perturbations, having been trained with domain randomization (random offsets in control) [oai_citation:23‡nature.com](https://www.nature.com/articles/s41467-024-52775-8?error=cookies_not_supported&code=7e8b5d76-2f55-4eae-96c0-2f080a6595a1#:~:text=Implementation%20of%20the%20control) [oai_citation:24‡pmc.ncbi.nlm.nih.gov](https://pmc.ncbi.nlm.nih.gov/articles/PMC11447118/#:~:text=This%20approach%20not%20only%20optimizes,to%20the%20real%20world%20experiment).  This work illustrates how an RL agent can effectively tune multi-parameter atomic experiments in situ.  In general, neutral-atom sensors provide high-dimensional continuous observations (images, currents) and allow smooth control; RL algorithms with continuous actions (e.g.\ DDPG or SAC) are appropriate.  For example, a neutral-atom interferometer (for rotation sensing or phase estimation) could use RL to adjust beam splitters or trap phases between measurement shots, based on prior outcomes, to maximize sensitivity.

\section{Reinforcement Learning Framework}
The RL agent interacts with the quantum sensor in episodic trials.  Each episode corresponds to a measurement protocol: the agent receives an {\it observation} (the state) describing the current sensor outputs and selects an {\it action} (control variables).  After the action, the quantum system evolves (possibly under an unknown parameter, such as a phase shift or magnetic field), and a measurement is performed to update the state.  The agent accumulates a {\it reward} that quantifies the information gained about the parameter.  Over many episodes, the agent optimizes its policy to maximize the expected cumulative reward.

\subsection{State, Action, and Reward}
We must define the state representation, action space, and reward to reflect our sensing goals.  {\bf State:} The state can include past measurement outcomes and any classical register of controls.  For example, in the double-dot setting one could use the measured current at various voltages (or pre-processed statistics) as the state vector [oai_citation:25‡nature.com](https://www.nature.com/articles/s41534-021-00434-x?error=cookies_not_supported&code=beb9a8cb-1d20-452c-ae4d-d852fc2325fa#:~:text=the%20efficient%20measurement%20of%20quantum,deep%20reinforcement%20learning%20for%20decision).  In a spin-based magnetometer, the state might include the previous spin measurement and the agent’s chosen pulse angle or amplitude.  In neutral-atom traps, the state could be images of the atom cloud or summary statistics (number, width).  Partial observability is inherent, so the agent’s state may consist of a history of recent observations (e.g.\ via an LSTM network) or a belief distribution from a Bayesian filter.

{\bf Actions:} Actions correspond to controllable experimental parameters.  For the quantum dot example, actions included moving to adjacent gate-voltage blocks [oai_citation:26‡nature.com](https://www.nature.com/articles/s41534-021-00434-x?error=cookies_not_supported&code=beb9a8cb-1d20-452c-ae4d-d852fc2325fa#:~:text=the%20efficient%20measurement%20of%20quantum,deep%20reinforcement%20learning%20for%20decision).  In general, actions can be discrete (e.g.\ change control setting to one of several levels) or continuous (e.g.\ adjust a laser detuning by a certain fraction).  In the spin magnetometer case, actions are continuous pulse parameters (field strength, direction, duration) [oai_citation:27‡arxiv.org](https://arxiv.org/html/2506.21475v1#:~:text=applying%20pulses%20of%20transverse%20fields,in%20quantum%20optimal%20control%20settings).  The agent selects a new control setting at each time step of the protocol.  The action space may be constrained by experimental limits (maximum pulse amplitude, etc.), which can be embedded as priors or constraints in a model-based RL algorithm [oai_citation:28‡arxiv.org](https://arxiv.org/abs/2501.14372#:~:text=and%20especially%20not%20for%20real,and).

{\bf Reward:} The reward should reflect the ultimate sensing objective.  A straightforward choice is the negative of the estimation error variance (or its logarithm), encouraging the agent to reduce uncertainty in the parameter.  For instance, if the goal is to estimate a phase $\phi$, one can compute the posterior variance after each measurement (using Bayesian update or via a particle filter) and assign the negative value as reward [oai_citation:29‡quantum-journal.org](https://quantum-journal.org/papers/q-2024-12-10-1555/#:~:text=especially%20given%20the%20capabilities%20of,adaptive%20and).  Alternatively, one can use the quantum Fisher information (QFI) directly: maximize QFI improves the Cram\'er--Rao bound [oai_citation:30‡arxiv.org](https://arxiv.org/html/2601.18953v1#:~:text=We%20have%20focused%20our%20description,based%20approach).  In practice, one might run a sequence of measurements in an episode and give a terminal reward equal to $-\log(\text{variance})$ at the end, or use incremental rewards related to information gain at each step.  The reward function may also include penalties for experimental cost or decoherence (e.g.\ long protocols may degrade signal).  Designing an informative reward is crucial, as it shapes the learning objective [oai_citation:31‡arxiv.org](https://arxiv.org/html/2601.18953v1#:~:text=literature%20is%20still%20smaller%20compared,transferred%20to%20quantum%20sensing%20tasks).

\subsection{RL Algorithms and Training}
Various RL algorithms can be employed, depending on the platform and problem dimensionality.  For discrete actions (e.g.\ choosing between several gate settings), value-based methods like Deep Q-Networks (DQN) are suitable; Nguyen \textit{et al.} used a DQN variant to scan QD voltage space [oai_citation:32‡nature.com](https://www.nature.com/articles/s41534-021-00434-x?error=cookies_not_supported&code=beb9a8cb-1d20-452c-ae4d-d852fc2325fa#:~:text=the%20efficient%20measurement%20of%20quantum,deep%20reinforcement%20learning%20for%20decision).  For continuous control (e.g.\ pulse amplitudes), policy-gradient or actor-critic methods like Deep Deterministic Policy Gradient (DDPG) or Soft Actor-Critic (SAC) are more natural [oai_citation:33‡pmc.ncbi.nlm.nih.gov](https://pmc.ncbi.nlm.nih.gov/articles/PMC11447118/#:~:text=illuminated%20with%20a%20resonant%20laser,final%20state%20of%20the%20atomic) [oai_citation:34‡arxiv.org](https://arxiv.org/html/2506.21475v1#:~:text=applying%20pulses%20of%20transverse%20fields,in%20quantum%20optimal%20control%20settings).  Off-policy methods (DQN, DDPG, SAC) are generally more sample-efficient, a key consideration since quantum measurements are costly [oai_citation:35‡arxiv.org](https://arxiv.org/html/2601.18953v1#:~:text=Figure%C2%A06%20%20shows%20a%20heuristic,space%2Ftabular%20problems).  In simulations, one can afford many episodes to train the agent.  In contrast, training directly on hardware may require domain randomization: varying unknown parameters during training so the agent generalizes, as done for the MOT by adding random perturbations to laser detuning [oai_citation:36‡nature.com](https://www.nature.com/articles/s41467-024-52775-8?error=cookies_not_supported&code=7e8b5d76-2f55-4eae-96c0-2f080a6595a1#:~:text=the%20MOT%2C%20and%20are%20typically,5%20G%20cm%E2%88%921) [oai_citation:37‡pmc.ncbi.nlm.nih.gov](https://pmc.ncbi.nlm.nih.gov/articles/PMC11447118/#:~:text=This%20approach%20not%20only%20optimizes,to%20the%20real%20world%20experiment).

If a high-fidelity simulator is available, model-aware RL (with automatic differentiation) can speed up learning by leveraging gradients through the known physics [oai_citation:38‡quantum-journal.org](https://quantum-journal.org/papers/q-2024-12-10-1555/#:~:text=especially%20given%20the%20capabilities%20of,adaptive%20and).  For example, if the Hamiltonian of the sensor is known, one can backpropagate through the simulation of the quantum evolution and measurement.  However, model-free approaches (not using gradients) are more general and simpler to implement, and have already shown success in sensing tasks [oai_citation:39‡arxiv.org](https://arxiv.org/html/2506.21475v1#:~:text=applying%20pulses%20of%20transverse%20fields,in%20quantum%20optimal%20control%20settings) [oai_citation:40‡nature.com](https://www.nature.com/articles/s41534-021-00434-x?error=cookies_not_supported&code=beb9a8cb-1d20-452c-ae4d-d852fc2325fa#:~:text=the%20efficient%20measurement%20of%20quantum,deep%20reinforcement%20learning%20for%20decision).

\section{Example: Phase Estimation with RL}
Consider an interferometric phase estimation task (e.g.\ optical interferometer or atomic Ramsey sequence).  The unknown phase $\phi$ is imprinted on a probe state, and one can apply a controllable phase shift $\theta$ before or after the probe.  A traditional adaptive strategy uses Bayesian feedback to choose $\theta$ based on prior outcomes.  An RL agent can play a similar role: the state includes the current Bayesian posterior (or measurement record), and the action is choosing the next phase $\theta$.  The reward can be related to the posterior variance in $\phi$ (e.g.\ via Fisher information).  By training the agent in simulation, it may discover non-intuitive adaptive protocols that approach Heisenberg scaling [oai_citation:41‡arxiv.org](https://arxiv.org/html/2601.18953v1#:~:text=In%20a%20nutshell%2C%20quantum%20metrology,While%20ground%20states%20can) [oai_citation:42‡quantum-journal.org](https://quantum-journal.org/papers/q-2024-12-10-1555/#:~:text=especially%20given%20the%20capabilities%20of,adaptive%20and).  For instance, model-aware RL could tune the probe state (squeezing parameter) and measurement basis in tandem to maximize QFI [oai_citation:43‡arxiv.org](https://arxiv.org/html/2601.18953v1#:~:text=We%20have%20focused%20our%20description,based%20approach).  Though explicit RL-based phase estimation beyond proof-of-principle is still being developed, the existing framework suggests it can match or surpass classical adaptive algorithms given sufficient training [oai_citation:44‡quantum-journal.org](https://quantum-journal.org/papers/q-2024-12-10-1555/#:~:text=especially%20given%20the%20capabilities%20of,adaptive%20and).

\section{Example: Magnetic Field Estimation with RL}
For magnetic field sensing, consider a spin-$1/2$ probe (e.g.\ an NV center or ion) in a background field $B_z$.  The agent controls a sequence of transverse control fields (pulses) to create an interferometric sequence.  After each evolution and measurement, the agent updates its strategy.  Cooke and Czischek modeled exactly this scenario: the agent learned a sequence of transverse pulses that maximizes sensitivity of a spin magnetometer, outperforming standard fixed strategies [oai_citation:45‡arxiv.org](https://arxiv.org/html/2506.21475v1#:~:text=applying%20pulses%20of%20transverse%20fields,in%20quantum%20optimal%20control%20settings).  They showed the agent generalizes to fields and parameters not seen in training, indicating robustness.  In a neutral-atom context, RL could optimize magnetic-field-induced phase shifts across an ensemble by adjusting trap or interaction times.  In any case, the reward can be defined as the negative variance of the estimated field, or equivalently as the (incremental) logarithm of the likelihood of the observed measurement given the true field.  The agent thus learns to probe the field distribution efficiently.

\section{RL Implementation Blueprint}
An actual implementation would follow these steps:
\begin{enumerate}
  \item {\bf Define the environment:} Choose the sensor platform and identify its controllable parameters and observables.  Build a simulation of the sensor dynamics (e.g.\ master equation or Monte Carlo) that includes the unknown parameter. If possible, validate the simulation against experiments.
  \item {\bf State representation:} Decide what information to feed to the agent.  This may be raw data (fluorescence images, spin readouts) or pre-processed features (means and variances). Include memory of previous steps if needed (e.g.\ via recurrent networks) to handle non-Markovian effects.
  \item {\bf Action space:} Specify the set of actions.  For continuous controls, discretization or parameterized actions can be used.  Ensure actions respect experimental limits (e.g.\ maximum pulse amplitude). 
  \item {\bf Reward function:} Choose a reward that captures precision.  A terminal reward $R = -\log(\text{estimation variance})$ is one option.  Alternatively, use stepwise rewards like Bayesian information gain.  Regularize the reward if necessary to avoid trivial strategies.
  \item {\bf RL algorithm:} Select an algorithm based on action type.  For example, use DQN for coarse adjustment of discrete parameters, or SAC/DDPG for smooth continuous control.  Set neural network architectures appropriate for the state (e.g.\ convolutional nets for images).
  \item {\bf Training:} Train the agent in simulation.  Introduce variability (randomize unknowns) to encourage generalization.  Use experience replay or batch updates to improve sample efficiency.  Monitor performance on validation scenarios.
  \item {\bf Deployment:} Transfer the trained policy to the real device.  If training was in-silico, consider fine-tuning on the hardware with a few adjustment episodes.  Include safeguards for safety (e.g.\ bounding pulses) and mechanisms to reset the environment.
\end{enumerate}

This outline treats the RL agent as a feedback controller that maximizes a precision-based objective.  The examples from recent experiments show each element is feasible: environments like OpenAI Gym were used for quantum dots, policy networks handled image-based state inputs in the MOT [oai_citation:46‡pmc.ncbi.nlm.nih.gov](https://pmc.ncbi.nlm.nih.gov/articles/PMC11447118/#:~:text=Here%2C%20we%20introduce%20deep%20RL,the%20code%2C%20running%20on%20a), and reward design (based on atom number or bias-triangle detection) successfully guided learning.  

\section{Conclusion}
Integrating reinforcement learning with quantum sensor platforms opens a path toward autonomous, high-precision estimation protocols.  We have described how an RL agent can interact with quantum dots, ion traps, or neutral-atom systems to optimize phase and field sensing.  Prior work provides evidence that RL can outperform heuristics: for instance, automatic RL-guided tuning of a double dot yielded faster characterization [oai_citation:47‡nature.com](https://www.nature.com/articles/s41534-021-00434-x?error=cookies_not_supported&code=beb9a8cb-1d20-452c-ae4d-d852fc2325fa#:~:text=the%20efficient%20measurement%20of%20quantum,deep%20reinforcement%20learning%20for%20decision), and RL-optimized control pulses improved magnetometer sensitivity [oai_citation:48‡arxiv.org](https://arxiv.org/html/2506.21475v1#:~:text=applying%20pulses%20of%20transverse%20fields,in%20quantum%20optimal%20control%20settings).  As hardware and simulators improve, RL-based quantum metrology promises fully data-driven adaptive sensors, potentially reaching quantum-limited precision.  Future directions include combining RL with quantum-enhanced probes (entangled states), multi-agent RL for distributed sensor networks, and hybrid classical-quantum RL where quantum processors simulate parts of the environment.  

\bibliographystyle{unsrt}
\bibliography{refs}
\end{document}
